% chapters/modelling.tex
\chapter{Numerical Modelling}
\label{ch:numerical-modelling}

% ---------- Chapter 4: Numerical Modelling ----------

\begin{figure}[htbp]
\centering

\staticsubfigure{figures/chapter4/fig4.1a.png}{}
\hfill
\staticsubfigure{figures/chapter4/fig4.1b.png}{}

\medskip

\staticsubfigure{figures/chapter4/fig4.1c.png}{}
\hfill
\staticsubfigure{figures/chapter4/fig4.1d.png}{}

\caption{The luminosity distance-redshift relation $d_L(z)$ (top) and distance modulus curve $\mu(z)$ ($\Omega_M$, $\Omega_X$) = (1, 0), solid; (0.05, 0), dotted; and (0.2, 0.8), dashed. (bottom) for various choices of the cosmological parameters $\Omega_M$ and $\Omega_X$ : ($\Omega_M$ , $\Omega_X$ ) = (1, 0) (solid) (0.05, 0) (dotted) and (0.2, 0.8) (dashed).}
\label{fig:4.1}
\end{figure}

\begin{figure}
    \centering
    \staticsubfigure{figures/chapter4/fig4.2a.png}{}
    \hfill
    \staticsubfigure{figures/chapter4/fig4.2b.png}{}
    \caption{An example of proper maximum step size in the Markov chain.  The walk (b) completely covers the parameter space in a random order, providing a smooth posterior (a).}
\label{fig:fig4.2}
\end{figure}


\begin{figure}
\centering

\staticsubfigure{figures/chapter4/fig4.3a.png}{}
\hfill
\staticsubfigure{figures/chapter4/fig4.3b.png}{}

\medskip

\staticsubfigure{figures/chapter4/fig4.3c.png}{}
\hfill
\staticsubfigure{figures/chapter4/fig4.3d.png}{}

\caption{Examples of improper maximum step sizes in the Markov chain. An incorrect step size causes the chain to inefficiently cover the parameter space: too small and the chain remains stuck in a small area, too large and the chain oscillates between maxima. The walks (d), (b) clearly show the problem, whereas the effect on the distributions (c), (a) is more subtle.}
\label{fig:4.3}
\end{figure}

The computational work involves four stages: writing appropriate code; the creation of SNe data; implementation of the Markov Chain Monte Carlo (\mcmc{}) method; creation of credible regions from the prior probabilities generated by \mcmc{}. The codes to do so are shown in \cref{tab:programs-written}. The programming language which I chose was \matlab{} due to its efficient handling of large arrays and the (relative) simplicity of generating visualisations of the data.\footnote{For a larger number of parameters or longer walks the low-level language Fortran is a more suitable program.}

\begin{table}[htbp]
\centering

\begin{threeparttable}
\caption{Programs written for the project.}
\label{tab:programs-written}

\begin{tabular}{@{} l l p{8cm} @{}}
\toprule
Filename & Type & Purpose \\
\midrule

\texttt{conf1} & Program &
Comparison of credible regions of single-parameter models for all data sets from Markov chains. \\

\texttt{conf2} & Program &
Visualisation of credible regions of 2 or 3 parameters for one data set from Markov chains. \\

\texttt{credibleregions} & Subfunction of \texttt{conf2} &
Calculation of credible regions for one parameter from Markov chains. \\

\texttt{friedmann} & Function &
Calculation of $\mu(z)$ from numerical integration of the Friedmann equations given a 5-tuple of parameters and an array of redshifts. \\

\texttt{friedmann\_rhs} & Subfunction of \texttt{friedmann} &
Array storing the ODEs to be solved. \\

\texttt{likelihood} & Function &
Global likelihood calculation from the $\mu(z)$ array generated by \texttt{friedmann}. \\

\texttt{montecarlo} & Program &
Implementation of the Metropolis--Hastings algorithm and generation of histograms and random walk plots. \\

\bottomrule
\end{tabular}

% Uncomment when table notes are needed.
% \begin{tablenotes}
% \footnotesize
% \item[a] Example table note.
% \end{tablenotes}
\end{threeparttable}
\end{table}


\section{Codes written for the project}
\label{sec:codes-written}

The codes required for the project varied widely in complexity (and consequently in debugging time). I wrote the codes starting with the clearest structure, then optimising them for speed. Two modifications caused the greatest increase in efficiency: firstly the replacement of for-loops with array operations; secondly minimising the memory used when writing data to files. Despite these measures, generating the Markov chains remained demanding in terms of computer time and hardware (\cref{sec:4.1.3}).

\subsection{Numerical solution of the Friedmann equations}
\label{sec:4.1.1}

The function designed to calculate $\mu(z)$ by solving the scaled Friedmann equations \cref{eq:2.7} was \texttt{friedmann}. The function calculated the initial conditions for the ODEs and passed them to the inbuilt \matlab{} ODE solver \texttt{ode45} along with the right hand side of the equations which were contained in the subfunction \texttt{friedmann\_rhs}. The output from the ODE solver was a 4 column array $\left[ \Omega_m (z_i), \Omega_X (z_i), H(z_i), D_L(z_i)\right]$, from which $D_L(z_i)$ was extracted to find $\mu(z_i)$ according to \cref{eq:1.13}.

The output of \texttt{friedmann} was the 2-column array of distance modulus versus redshift $[z_i \,, \mu(z_i)]$. Plotting this generated a redshift-distance modulus curve. I compared my own results \cref{fig:fig4.1b} and \cref{fig:fig4.1d} to those \cref{fig:fig4.1a} and \cref{fig:fig4.1c} published by Hogg in \cite{Hogg2000} to ensure that \texttt{friedmann} produced correct $d_L(z)$ and $\mu(z)$ values. Due to the assumption of a flat spatial geometry (with $\Omega_m + \Omega_X = 1$) I could not produce the curve for $(\Omega_m = 0.05, \, \Omega_X = 0)$. Therefore I was able to verify that the script both correctly solved the Friedmann equations and correctly converted from the distance-redshift relation to the observed quantity given in the data.

\subsection{Likelihood calculations}
\label{sec:4.1.2}

The random walk requires the calculation of the likelihood $L(x)$ at each selected point $x$ in the parameter space. The point with the largest likelihood corresponds to the set of parameter values which are most likely to have generated the $\mu(z)$ curve formed from the data.   The analytical value \cref{eq:4.1} for $L(x)$ is the product of the likelihood according to each point $\mu(z_j)$ in the data. Since the error in $\mu$ is well-characterised by a Gaussian distribution, the individual likelihoods are \cref{eq:4.1} where at the redshift $z_j$ , $\sigma$ is the standard deviation of the uncertainty in the distance modulus measurement $\mu(z_j)$ (this was provided explicitly in the data) and $\mu(x)$ is the distance modulus calculated by the model with parameter 5-tuple $x$. The individual likelihood is maximised for the curve which, at each $z_j$ in the chosen SNe data set, matches as closely to the distance modulus $\mu_j$ of the SN there, i.e. when $\mu(x) - \mu(z_j) = 0$.

% Equation (4.1)
\begin{equation}
    \mathcal{L}(x)
    = \prod_{j} \mathcal{L}_{j}(x)
    \quad
    \text{where}
    \quad
    \mathcal{L}_{j}(x)
    = \mathcal{L} \left( \mu{} ( z_{j};x ) \right)
    = \frac{1}{\sqrt{2\pi{}\sigma{}}} \exp 
    \left( 
        - \frac{\left( \mu{} (x) - \mu{} (z_{j}) \right)^{2]}}{2 \sigma{}^{2} }
    \right)
\label{eq:4.1}
\end{equation}

Since the global likelihood is the product of Gaussians, computational accuracy is maximised by calculating the logarithms of the individual values to avoid overflow. To perform global likelihood calculations for each set of parameter values I wrote a subfunction called likelihood rather than using a pre-written \matlab{} routine. The scalar output from likelihood was passed to the main program which used this and the previous value to determine the next model to pick in parameter space as described in \cref{alg:2.1}.

\subsection{The Metropolis-Hastings algorithm}
\label{sec:4.1.3}

The Metropolis-Hastings algorithm (\cref{alg:2.1}) is a random walk which produces Markov chains that accept or reject a proposed step on the basis of the likelihood ratio between the steps. The main program \texttt{montecarlo} looped over each maximum step size, performing the Metropolis-Hastings algorithm for a given walk length, writing the parameter values and likelihood to file at each step, before creating a histogram and a line graph for each of the varying parameters in the walk. The algorithm contains three computationally important factors: the number of steps in the walk, the dimensionality of the parameter space and the size of the neighbourhood around each point. As discussed in \cref{sec:4.3}, although writing the code is straightforward, ensuring that it truly approximates the posteriors from the longer crude Monte Carlo method is not so simple.


\section{Generation of the projected WFIRST set}
\label{sec:4.2}

The projected \textsc{wfirst} data is designed to illustrate improvements in constraints on the parameters given a larger Ia SNe data set from a single survey. The first step in making the projected data was to generate 2000 redshifts from a random distribution $0.1 < z < 1.7$ using \texttt{randn}, a built-in \matlab{} (pseudo-)random number generator. The Friedmann integrator (\cref{sec:4.1.1} calculated the distance moduli at each redshift. 

Two types of errors were now introduced: those due to the instrument and those from uncontrollable sources. The first arises from the photometric calibration and is estimated by the Dark Energy Task Force report in \cite{Albrecht2006} to be a linear function of redshift $k(1 + z)$, with $k$ a constant between $0.01$ for an \lq\lq{}optimistic\rq\rq{} view and $0.05$ for a \lq\lq{}pessimistic\rq\rq{} one. Since other errors are not quantified by the working group, I estimated them from the most recent SNe survey, the Union 2 compilation \cite{Amanullah2010}. These additional errors (potentially arising from a variety of sources discussed in \cref{sec:3.1} and \cite{Leibundgut2009,Lidman2010b}) were assigned by random selection from the Union2 set to the artificial SNe. The total uncertainty in the distance moduli is then the sum of the systematic error and the random error estimates. Each data point was then displaced from the concordance model curve by the error bar size multiplied by a random number from $[-1, 1]$. Two sets were made: the snap (snap 02) set using the pessimistic (optimistic) error from the JDEM report and the same magnitude ($0.2x$ magnitude) errors as the Union2 survey, maintaining the scaling of the quantified errors.

\section{Simulation of the posterior distributions}
\label{sec:4.3}

The posterior distributions for each Ia SNe compilation and model are approximated by histograms created using the Metropolis-Hastings version of Markov chain Monte Carlo. The requirements for a reasonable approximation to the true (unknown) posterior are:

\begin{itemize}
    \item Sufficient burn-in time for the chain to reach an equilibrium.
    \item A suitable maximum step size to explore the parameter space.
    \item Initial conditions (a \lq\lq{}zero-parameter model\rq\rq{}) which do not introduce bias into the walks.
\end{itemize}

Ensuring that all these conditions were met without resorting to highly complex code and parallel computing was a non-trivial task involving much trial and error. This constituted the majority of the project: although simulations were started in May, only 4065 hours of CPU time actually produced data suitable for further processing.\footnote{The output files were timestamped so that it was possible to tell the duration of the walk.} The rest of the time was spent discovering the aforementioned problems in practice, rewriting the code and running new simulations to determine whether the issue had been resolved. 

In this section, I examine the problems faced in more detail and outline the methods used to ameliorate them. The convergence length of the Markov chain cannot be predetermined. Although the \mcmc{} formalism guarantees convergence for most PDFs, it places no bounds on the number of iterations required before convergence is achieved \cite{Gregory2005,Dunkley2005,Christensen2001}. Existing methods are empirical: one runs a series of tests for non-convergence, the negative results of which imply---but do not confirm---that the chain has converged \cite{Christensen2001}. These methods, which include using second-order Markov chains to calculate a bound on the accuracy of the estimations and comparing the values from multiple chains \cite{Dunkley2005}, were considered unnecessary for such a small number of parameters (compared to eleven or more in \cmb{} calculations \cite{LewisBridle2002}). Instead, following \cite{Gregory2005}, the burn-in time preceding convergence was estimated at $1\%$ of the total walk length, provided that the posterior produced was smooth. 

After some experimenting to find a smooth histogram of $\Omega_m$ from the Constitution compilation, the initial walk length estimate of $5 \times 10^4$ steps was increased to $1 \times 10^5$ for the smaller-error observational data (Constitution 0.2, Gold 0.2, Union 0.2) and $2 \times 10^5$ for the projected sets (Snap, Snap 0.2) for one-parameter models. The number of steps required to produce smooth posteriors scales with the dimensionality of the parameter space: (approximately) linearly for \mcmc{} methods \cite{Dunkley2005} compared to exponentially in the crude MC scenario \cite{Christensen2001}. In practice, this linear scaling is not achieved \cite{Christensen2001,Dunkley2005,LewisBridle2002}. Therefore a balance between smoothness and running time was estimated by multiplying the one-parameter walk length tenfold for each additionally varying parameter. The exception to this was multiple-parameter models involving the time-varying dark energy equation of state $w_1$, which had a long burn-in time: these walks had to be redone with twice the existing steps for smoothness. Computational efficiency required that the total number of steps in the chain was both sufficiently large to widely explore the parameter space (\cref{fig:fig4.2b}) and sufficiently small that the program runs in a sensible amount of time given the time span available for the project. 

The second problem in the algorithm is the possibility of sampling bias introduced by the the neighbourhood size. This is governed by a maximum step size: the actual steps will be a fraction of this (depending on the random multiplier chosen). If the maximum step is too small, the walk does not leave a small subset of the parameter volume (\cref{fig:fig4.3b}). Consequently, it will show fine detail (\cref{fig:fig4.3a}), but only in a small section of the characteristic width, leading to an incorrect posterior. Conversely, if the step size is too large, the chain (\cref{fig:fig4.3d}) will oscillate between regions of high likelihood (\cref{fig:fig4.3c}) without exploring the parameter space in-between. This results in a posterior which sparsely outlines peaks in the likelihood, but does not give any accurate shape to the rest of the space. Thus the same initial conditions can produce two very different chains if different maximum step sizes are selected. Although a method called parallel tempering \cref{sec:6.1} exists which efficiently covers the parameter space using information from a variety of step sizes, it was too complex to implement in this project. Consequently the walk had to be repeated for each model over a range of different step sizes incremented by factors of two using a for-loop. It became evident that there was no step size which was suitable for all parameters, e.g. the Hubble parameter $h$ required a small step because of the small characteristic width of its prior, whereas the dark energy equation of state parameters $w_0$ (constant) and $w_1$ (time-varying) produced smooth walks with much larger step sizes. The core step sizes were $0.1 \times 2k ,\, k \in {-2, -1, 0, 1, 2}$ and additional sizes were added after the walks from these were examined: if the $0.025$ step was too large then one of $0.0125$ was run; if the $0.4$ step was too small, steps of $0.8$ and even $1.6$ were run.

The final issue to address is the choice of initial conditions for the walk. If the initial conditions were too far from the likely parameter values then a large number of steps would be wasted as burn-in time: the section of the walk in which the chain travels towards high-likelihood regions \cite{Gregory2005}. Consequently the initial conditions chosen corresponded to the concordance ($\Lambda$CDM) model, an accepted estimate to the correct values of the cosmological parameters \cite{Hartle2003}. The values of the concordance model vary slightly from reference to reference: I adopted the 1-significant figure values used in \cite{Hartle2003}: $\Omega_m = 0.3; w_0 = \Lambda = -1$; with the addition of $w_1 = \gamma = 0$ (by definition the concordance model has neither time-variation in the dark energy equation of state nor interaction). I chose to write the Hubble constant parameter as $h/h_{72}$ so that a value of 1 gave $H_0 = 72\,\mathrm{km}\,\mathrm{s}^{-1}\,\mathrm{Mpc}^{-1}$ , the current best estimate for $H_0$ \cite{Riess2009}.

Once all three criteria were met, the Markov chains produced use approximations to the true posterior PDFs. This enabled the smooth walks with suitable step sizes (\cref{fig:fig4.2}) to be selected for the next stage in the project: finding credible regions.


\section{Credible regions from posteriors}
\label{sec:4.4}

Not all of the histograms from the Markov chains were reasonable approximations to the posterior PDFs for that model. I categorised all successful Markov chain histograms:
\begin{equation*}
    (\text{8 data sets})
    \times
    \left[
        \binom{4}{2} 
        +
        \binom{4}{1}
    \right] \text{models}\footnote{The zero-parameter model is $\Lambda$CDM.}
    \times
    (\text{6 step sizes})
    =
    8 \times (6+4) \times 6
    =
    480\ \text{histograms}.
\end{equation*}
as unsuitable (maximum step size either too small or too large) or suitable (satisfying all three criteria). I then selected the smoothest histograms (if there were multiple maximum step sizes which produced a good result) for further analysis. The next step in the analysis was the creation of credible regions for 1d and 2d subspaces of the 5d parameter space.

\subsection{Credible regions for individual parameters}
\label{sec:4.4.1}

Although we desire a single estimate for each parameter, the posteriors are not delta functions, so there is an inherent width to the high-probability region whence the \lq\lq{}best fit\rq\rq{} estimate is selected. To examine the posterior PDFs in greater detail, I wrote the graphing program conf1. This read in the Markov chain data from a set of files and output line graphs of each posterior. The process involved binning the data, ordering the bins by population, then selecting those regions within which lay $68\%$, $95\%$ and $99\%$ of the points in the chain. These regions were plotted as line graphs shaded according to increasing probability content. The mode and standard deviation were also plotted as a point with errorbars above the curve. Thus I was able to facilitate comparison between the different data sets.

\subsection{Credible regions for multiple parameters}
\label{sec:4.4.2}

Models with more than one free parameter can be examined in a similar manner to those with only a single free parameter. Although the theory is the same, the binning is based upon the values of all parameters simultaneously. Therefore I needed to write not only a program to generate contour maps of the credible regions (\texttt{conf2}) but also one to calculate them (\texttt{credibleregions}). Since the Metropolis- Hastings algorithm moves in all directions within the 2d parameter space, the value of both parameters is related. Unlike the existing \matlab{} function \texttt{histc}, which bins the parameter values in the random walk independently of one another, \texttt{credibleregions} bins the 2d subspace according to each pair of values. This retained any patterns in the walk, allowing me to look for correlations between parameters.
